\newcommand\UniversiteAdi{Niğde Ömer Halisdemir Üniversitesi}
\newcommand\BolumAdi{MEKATRONİK BÖLÜMÜ}
\newcommand\DersKodu{MKT2002}
\newcommand\DersAdi{BİLGİSAYARLI KONTROL SİSTEMLERİ}
\newcommand\SinavAdi{Finale Hazırlık}
\newcommand\SinavTarihi{}
\newcommand\SinavSaati{}
\newcommand\SinavSuresi{}

\pagestyle{fancy}
\fancyhf{} % clear existing header/footer entries
\fancyhead[R]{Öğrenci No:\hspace{4.5cm}}
\fancyhead[L]{Ad Soyad:\hspace{7cm}}
\noindent \includegraphics[width=0.1\textwidth]{logo}
\begin{tabular}{
    p{0.15\linewidth}
    p{0.15\linewidth}
    p{0.2\linewidth}
    p{0.1\linewidth}
    p{0.15\linewidth}}
    \multicolumn{5}{c}{\textbf{\BolumAdi}}\\
    \multicolumn{5}{c}{\textbf{\DersAdi}}\\\hline
    \multicolumn{1}{|r|}{Ders Kodu:}&
    \multicolumn{1}{|c|}{\DersKodu}&
    \multicolumn{1}{|c|}{}& 
    \multicolumn{1}{|r|}{Tarih:}&
    \multicolumn{1}{|c|}{\SinavTarihi} \\\hline
    \multicolumn{1}{|r|}{Sınav Türü:}&
    \multicolumn{1}{|c|}{\SinavAdi}&  
    \multicolumn{1}{|c|}{}&
    \multicolumn{1}{|r|}{}&
    \multicolumn{1}{|c|}{\SinavSaati}\\\hline
    \multicolumn{1}{|r|}{Dönemi:}&
    \multicolumn{1}{|c|}{2024-2025}&
    \multicolumn{1}{|c|}{}&
    \multicolumn{1}{|r|}{}&
    \multicolumn{1}{|c|}{\SinavSuresi} \\\hline
    &&&&\\
\end{tabular}\\\\
\noindent\textbf{Soru:} Bir açık çevrim transfer fonksiyon modeli
\begin{equation*}
    G(s)=\frac{1}{s^2+s+1}
\end{equation*}
olarak verilmiştir. Kapalı çevrim için yerleşme zamanı $t_s=1\,s$ ve aşım $\%10$ isterleri verilmiştir.

Aşım isteri ile 
\begin{equation}
\begin{split}
    \zeta&=-\frac{\log(os)}{\sqrt{\pi^2+\log{os}^2}}\\
    &=0.591
\end{split}
\end{equation}
ve yerleşme zamanı isteri ile 
\begin{equation}
\begin{split}
    \omega_n&=\frac{4}{\zeta t_s}\\
    &=6.7659
\end{split}
\end{equation}
elde edilir. İsterleri sağlayacak polinom
\begin{equation}
\begin{split}
    p(s)&=s^2+2\zeta \omega_n s+\omega_n^2\\
    &=s^2+8s+45.7844\\
\end{split}
\end{equation}
olarak önerilir.

\noindent\textbf{a)} P kontrolör tasarlayınız.
P kontrolör
\begin{equation}
    F(s)=k
\end{equation}
olmak üzere kapalı çevrim transfer fonksiyonu 
\begin{equation}
\begin{split}
    T(s)&=\frac{F(s)G(s)}{1+F(s)G(s)}\\
    &=\frac{k\frac{1}{s^2+s+1}}{1+k\frac{1}{s^2+s+1}}\\
    &=\frac{k}{s^2+s+1+k}
\end{split}
\end{equation}
olarak elde edilmektedir. Tasarım için 
\begin{equation}
\begin{split}
    1&=8\\
    1+k&=45.7844
\end{split}
\end{equation}
eşitliği çözülmelidir. Görüldüğü üzere çözüm mevcut değildir. Sadece aşım için çözülmek istenirse, \begin{equation}
\begin{split}
    p(s)&=s^2+2\zeta \omega_n s+\omega_n^2\\
    &=s^2+1.1824\omega_n s+\omega_n^2
\end{split}
\end{equation}
kullanılarak 
\begin{equation}
\begin{split}
    1.1824\omega_n&=8\\
    1+k&=\omega_n^2
\end{split}
\end{equation}
probleminin çözümü $\omega_n=0.8458$ ve $k=-0.2846$ elde edilmektedir. Kapalı çevrim transfer fonksiyonu
\begin{equation}
    T(s)=\frac{-0.2846}{s^2 + s + 0.7154}
\end{equation}
olarak hesaplanmaktadır. Basamak yanıtı Şekil~\ref{fig:p_kontrol} ile gösterilmektedir.
\begin{figure}[!htb]
    \centering
    \includegraphics[width=0.75\textwidth]{p_kontrol}
    \caption{P kontrolör}
    \label{fig:p_kontrol}
\end{figure}
\clearpage

\noindent\textbf{b)} PD kontrolör tasarlayınız.
PD kontrolör
\begin{equation}
    F(s)=k_d s+k_p
\end{equation}
olmak üzere kapalı çevrim transfer fonksiyonu 
\begin{equation}
\begin{split}
    T(s)&=\frac{F(s)G(s)}{1+F(s)G(s)}\\
    &=\frac{(k_d s+k_p)\frac{1}{s^2+s+1}}{1+(k_d s+k_p)\frac{1}{s^2+s+1}}\\
    &=\frac{(k_d s+k_p)}{s^2+(k_d+1)s+(k_p+1)}
\end{split}
\end{equation}
olarak elde edilmektedir. Tasarım için 
\begin{equation}
\begin{split}
    k_d+1&=8\\
    k_p+1&=45.7844
\end{split}
\end{equation}
eşitliğinin çözümü $k_d=7$ ve $k_p=44.7844$ elde edilmektedir. Kapalı çevrim transfer fonksiyonu
\begin{equation}
    T(s)=\frac{7 s + 44.78}{s^2 + 8 s + 45.78}
\end{equation}
olarak hesaplanmaktadır. Basamak yanıtı Şekil~\ref{fig:pd_kontrol} ile gösterilmektedir.
\begin{figure}[!htb]
    \centering
    \includegraphics[width=0.75\textwidth]{pd_kontrol}
    \caption{PD kontrolör}
    \label{fig:pd_kontrol}
\end{figure}
\clearpage

\noindent\textbf{c)} PI kontrolör tasarlayınız.
PI kontrolör
\begin{equation}
    F(s)=k_p+\frac{k_i}{s}
\end{equation}
olmak üzere kapalı çevrim transfer fonksiyonu 
\begin{equation}
\begin{split}
    T(s)&=\frac{F(s)G(s)}{1+F(s)G(s)}\\
    &=\frac{\frac{k_ps+k_i}{s}\frac{1}{s^2+s+1}}{1+\frac{k_ps+k_i}{s}\frac{1}{s^2+s+1}}\\
    &=\frac{k_ps+k_i}{s^3+s^2+(k_p+1)s+k_i}
\end{split}
\end{equation}
olarak elde edilmektedir. 3. dereceden bir kapalı çevrim karakteristik polinom elde edilmesi sebebiyle 
isterleri sağlayacak polinom
\begin{equation}
\begin{split}
    p(s)&=(s^2+2\zeta\omega_n s+\omega_n^2)(s+p)\\
    &=(s^2+8s+45.7844)(s+p)\\
    &=s(s^2+8s+45.7844)+p(s^2+8s+45.7844)\\
    &=s^3+8s^2+45.7844s+ps^2+8ps+45.7844p\\
    &=s^3+(p+8)s^2+(8p+45.7844)s+45.7844p
\end{split}
\end{equation}
kullanılmalıdır.
Tasarım için problem
\begin{equation}
\begin{split}
    1&=p+8\\
    k_p+1&=8p+45.7844\\
    k_i&=45.7844p
\end{split}
\end{equation}
ve çözümü $k_i=-320.49$, $k_p=-11.21$, $p=-7$ şeklinde elde edilmektedir. $p$'nin negatif çıkması çözümün çalışmadığını göstermektedir. 
\clearpage

\noindent\textbf{d)} PID kontrolör tasarlayınız.
PID kontrolör
\begin{equation}
    F(s)=k_p+\frac{k_i}{s}+k_d s
\end{equation}
olmak üzere kapalı çevrim transfer fonksiyonu 
\begin{equation}
\begin{split}
    T(s)&=\frac{F(s)G(s)}{1+F(s)G(s)}\\
    &=\frac{\frac{k_ds^2+k_ps+k_i}{s}\frac{1}{s^2+s+1}}{1+\frac{k_ds^2+k_ps+k_i}{s}\frac{1}{s^2+s+1}}\\
    &=\frac{k_ds^2+k_ps+k_i}{s^3+(k_d+1)s^2+(k_p+1)s+k_i}
\end{split}
\end{equation}
olarak elde edilmektedir. 3. dereceden bir kapalı çevrim karakteristik polinom elde edilmesi sebebiyle 
isterleri sağlayacak polinom
\begin{equation}
\begin{split}
    p(s)&=(s^2+2\zeta\omega_n s+\omega_n^2)(s+p)\\
    &=(s^2+8s+45.7844)(s+p)\\
    &=s(s^2+8s+45.7844)+p(s^2+8s+45.7844)\\
    &=s^3+8s^2+45.7844s+ps^2+8ps+45.7844p\\
    &=s^3+(p+8)s^2+(8p+45.7844)s+45.7844p
\end{split}
\end{equation}
kullanılmalıdır.
Tasarım için problem
\begin{equation}
\begin{split}
    k_d+1&=p+8\\
    k_p+1&=8p+45.7844\\
    k_i&=45.7844p
\end{split}
\end{equation}
olarak tanımlanmaktadır. 4 bilinmeyen 3 denklem için sonsuz farklı çözüm mevcuttur. Parametrik çözüm 
\begin{equation}
\begin{split}
    k_d&=p+7\\
    k_p&=8p+44.7844\\
    k_i&=45.7844p
\end{split}
\end{equation}
olarak elde edilir ve $p$ parametresi pozitif seçilebilir. Kapalı çevrim transfer fonksiyonu 
\begin{equation}
\begin{split}
    T(s)&=\frac{(p+7)s^2+(8p+44.7844)s+45.7844p}{s^3+(p+8)s^2+(8p+45.7844)s+45.7844p}
\end{split}
\end{equation}
olarak elde edilmektedir. İsterlerin sağlanması için seçilen kapalı çevrim transfer fonksiyonu
\begin{equation}
\begin{split}
    T(s)&=\frac{45.7844p}{s^3+(p+8)s^2+(8p+45.7844)s+45.7844p}
\end{split}
\end{equation}
ve $p\geq 5\zeta\omega_n$ olmalıdır. Fakat, pay ifadesi iki adet sıfır içermektedir ve bu sıfırların da 5 kat uzakta olması gerekmektedir. Hem 3. kutbu hem de sıfırları en uzağa yerleştirecek $p$ değeri tarama yöntemi ile
\begin{equation}
\begin{split}
    p&=19.9\\
    k_p&=203.984\\
    k_d&=26.9\\
    k_i&=911.11
\end{split}
\end{equation}
ve kapalı çevrim transfer fonksiyonu
\begin{equation}
    T(s)=\frac{26.9 s^2 + 204 s + 911.1}{s^3 + 27.9 s^2 + 205 s + 911.1}
\end{equation}
olarak hesaplanmıştır. Basamak yanıtı Şekil~\ref{fig:pid_kontrol} ile gösterilmektedir.
\begin{figure}[!htb]
    \centering
    \includegraphics[width=0.75\textwidth]{pid_kontrol}
    \caption{PID kontrolör}
    \label{fig:pid_kontrol}
\end{figure}
